% -*- coding: utf-8 -*-

\begin{zhaiyao}
\begin{spacing}{1.5}
学位论文是本科生工作成果的集中体现，是本科生申请学士学位的主要依据，也是社会重要的文献资料。为了进一步推进我校本科生学位论文的规范化，提高写作质量，我们编写了《南开大学本科生学位论文写作规范（试行）》，供申请学位的本科生参考。

\textcolor{red}{注意：中文摘要应简要说明毕业论文所研究的内容、目的、方法、结论、主要成果和特色，字数一般在 200 至 300 之间，语言力求精炼。}
\end{spacing}

\vspace{\baselineskip}
\noindent{\zihaoxiaosi\bfseries 关键词：}

\hspace*{2em}软件工程；学位论文；写作规范

\textcolor{red}{注意：论文的关键词应是从其题目、层次标题和正文中选出来的，能反映论文主题概念的词或词组，一般为 3—5 个。}
\end{zhaiyao}

\begin{abstract}
\begin{spacing}{1.5}
The dissertation is concentrated expression of undergraduate work, which is not only the main basis for undergraduate applicants for a bachelor's degree, but also a socially important document.

\textcolor{red}{Abstract 内容与中文摘要相对应。}

\textcolor{red}{注意：不要用“google 翻译”、“金山快译”、“有道翻译”等软件翻译后直接 copy，一定要认真修改，确保语法、句法没有问题，同时尽可能合乎英语常用语规范。}

（注明：字体要求：小四字号、Times New Roman）
\end{spacing}

\vspace{\baselineskip}
\noindent{\zihaoxiaosi\bfseries Keyword}

\hspace*{2em}Software Engineering; Dissertation; Writing specifications
\end{abstract}
